% Lösungen Einheit 4 von 4 – Periodische Vorgänge modellieren (zusammenbau v0.1)
\einheitenkopf{Einheit 4 von 4 · Periodische Vorgänge modellieren}

\erg{21}{a) wiederholt sich \quad b) periodisch – nach jeder Hin- und Herbewegung beginnt dieselbe Bewegung von vorn \quad c) 6 s (hell und dunkel zusammen) \quad d) 27; 13 (in °C) \quad e) Mittellinie = 20 °C; Amplitude = 7 °C \quad f) a = 7; b = 360 : 24 = 15; d = 20; also y = 7 · sin(15 · t) + 20 \quad g) y = 16 · sin(45 · t) + 18 \quad h) y ≈ 26,1 – um 12 Uhr ist es im Gewächshaus etwa 26 °C warm \quad i) t = 6, also um 14 Uhr (dort ist 15 · t ein rechter Winkel) \quad j) ≈ 29 m \quad k) nur ungefähr: jeder Tag ist anders warm, bei Regen fehlt der Höchstwert – das Modell liefert Näherungswerte, keine Messwerte \quad l) wiederholt sich: nur die Sinuskurve; gleiche Zunahme je Schritt: nur die Gerade \quad m) B – A ist eine Gerade, aber 3 \% von einem wachsenden Betrag sind jedes Jahr mehr, der Graph muss gekrümmt sein; C beginnt bei 0 €, das Guthaben beginnt aber bei 500 € \quad n) a) (1); b) A; c) y ≈ 16,3 – rund 100 Tage nach dem 21. März, Ende Juni, ist der Tag gut 16 Stunden lang}
\erg{22}{a) Fehler: Graph als Bild des Vorgangs gedeutet. Der Graph zeigt die Höhe über der Zeit; die Gondel selbst fährt auf einem Kreis. \quad b) Es wächst immer weiter und kehrt nie zu einem früheren Wert zurück – nichts wiederholt sich.}
